\begin{table*}[t]
\caption{\bt{RelBench: predictive quality, valid solutions, and reuse cost. Each method runs three tasks at five cutoffs per database.}{RelBench：预测质量、有效方案与复用代价。每种方法在每个数据库上运行三个任务、五个截点。}}
\label{tab:relbench}
\centering
\ifbilingual\footnotesize\else\small\fi
\setlength{\tabcolsep}{2.8pt}
\renewcommand{\arraystretch}{1.08}
\newcommand{\RbPanelHeader}[4]{%
  \multicolumn{12}{@{}l}{\textbf{#1}}\\[2pt]
  & \multicolumn{6}{c}{\bh{Reuse cost: cutoffs 2--5}{复用代价：截点 2--5}}
  & \multicolumn{3}{c}{\bh{Mean test score: cutoffs 1--5}{平均测试分数：截点 1--5}}
  & \multicolumn{2}{c}{\bh{Quality summary}{质量汇总}}\\
  \cmidrule(lr){2-7}\cmidrule(lr){8-10}\cmidrule(l){11-12}
  \bh{Method}{方法} & \bh{Tokens\\(M)}{token\\（百万）} & \bh{Token\\ratio}{token\\比率}
  & \bh{LLM\\calls}{LLM\\调用} & \bh{Tool\\calls}{工具\\调用}
  & \bh{Obs.\\min}{观测\\分钟} & \bh{Cost\\(USD)}{美元}
  & #2 & #3 & #4 & \bh{Valid\\sessions}{有效会话} & \bh{Quality\\tasks}{达标任务}\\
  \midrule
}
\begin{tabular*}{\textwidth}{@{\extracolsep{\fill}}lrrrrrrrrrcc@{}}
\toprule
\RbTableBody
\bottomrule
\end{tabular*}
\par\smallskip
\parbox{\textwidth}{\footnotesize\raggedright
\bt{Cost columns total 12 sessions at cutoffs 2--5, including memory operations; token ratios are relative to no memory. Test scores average all five cutoffs, using the entity-table reference for sessions without a valid solution. \emph{Valid sessions}: non-fallback solutions (/15). \emph{Quality tasks}: tasks retaining no-memory predictive quality (/3), with $\mathrm{AUC}\geq\mathrm{AUC}_0-0.01$ or $\mathrm{MAE}\leq 1.02\,\mathrm{MAE}_0$. Subscript 0 denotes no memory; best task means are bold. Time and USD accounting are specified in Section~\ref{sec:relbench-eval}.}{代价列合计截点 2--5 的 12 个会话，含记忆操作；token 比率以无记忆为参照。测试分数取五个截点的平均，无有效方案的会话采用实体表参照分数。有效会话计数排除参照回退（共 15 个）；达标任务指保持无记忆预测质量的任务数（共 3 项），满足 $\mathrm{AUC}\geq\mathrm{AUC}_0-0.01$ 或 $\mathrm{MAE}\leq 1.02\,\mathrm{MAE}_0$。下标 0 表示无记忆；各任务最佳平均分加粗。时间与美元计量见第~\ref{sec:relbench-eval} 节。}
}
\end{table*}
